\section{The delay test, and what the archimedean expansion selects}
\label{sec:delay}

Condition~\ref{cond:bilinear} is the one test in this paper a theory can fail.
This section runs it. The answer is that the test as it was first posed is failed
by every correlator and passed by a free Hamiltonian, so it discriminates nothing;
but the \emph{expansion} of the delay does discriminate, one coefficient at a
time, and it recovers the arithmetic invariants of an $L$-function in order.

\subsection{The object the test can be run on}

The natural reading of Condition~\ref{cond:bilinear} is to take a protected
two-point function on the open Wilson line of \cite{OWL} and extract the phase of
its spectral kernel. There is no such correlator: \cite{OWL} establishes the
existence of supersymmetric open Wilson lines --- only in the superconformal
versions of the $\mathcal N=2$ theories, and in $\mathcal N=4$ with a defect
hypermultiplet, where the configuration is a semicircle conformally related to the
ray --- and computes \emph{expectation values}. It contains no two-point function
of open Wilson lines and no correlator of insertions along the line.\footnote{Its
abstract was read directly; the body through a secondary reading, so this is
"not found" rather than "certainly absent".}

The test is therefore run on the object that does exist and is known exactly, the
defect CFT on the $\frac12$-BPS Wilson line, whose protected two-point functions
are $C_\Delta(\lambda)/|t_{12}|^{2\Delta}$ with $\Delta=1$ for the five scalars
\cite{GRT}, $\frac32$ for the fermions, $2$ for the displacement
\cite{CHMS}, and $L$ for the $(Y\!\cdot\!\Phi)^L$ tower \cite{GiombiKomatsu} ---
every dimension protected, the coupling dependence entirely in the normalization.

\subsection{The delay of a defect two-point function}

In the ray coordinate of Section~\ref{subsec:conformal}, $x=\log r$ and
$u=x_1-x_2$, a defect two-point function is
$G_\Delta(u)=(2\sinh\frac u2)^{-2\Delta}$, the conformally covariant form of
$C/|r_1-r_2|^{2\Delta}$; the normalization is a positive constant and drops out of
every phase. Its causal symbol is
\begin{equation}\label{eq:defectsymbol}
 \widehat k_+(\tau)=\int_0^\infty e^{-i\tau u}G_\Delta(u)\dd u
 =\Gamma(1-2\Delta)\,\frac{\Gamma(\Delta+i\tau)}{\Gamma(1-\Delta+i\tau)} ,
\end{equation}
checked against quadrature. Using
$\frac{d}{d\tau}\arg\Gamma(a+i\tau)=\re\psi(a+i\tau)$ and
$\psi(1-z)-\psi(z)=\pi\cot\pi z$:

\begin{proposition}[The delay of a defect primary]\label{prop:defectdelay}
For every $\Delta$,
\[
 \mathcal T_\Delta(\tau)=-\frac{d}{d\tau}\arg\widehat k_+(\tau)
 =\pi\re\cot\big(\pi(\Delta+i\tau)\big)
 =\frac{\pi\sin(2\pi\Delta)}{\cosh(2\pi\tau)-\cos(2\pi\Delta)} .
\]
It decays like $2\pi\sin(2\pi\Delta)e^{-2\pi\tau}$; for $0<\Delta<\frac12$ its
total is $\int_0^\infty\mathcal T_\Delta=\pi(\frac12-\Delta)$; and it is periodic
in $\Delta$ with period one.
\end{proposition}

\begin{corollary}[The protected sector gives exactly nothing]\label{cor:protected}
$\mathcal T_\Delta\equiv0$ whenever $2\Delta\in\Z$, hence for every protected
operator of the $\frac12$-BPS Wilson-line defect CFT. Not small: zero, at every
frequency and every coupling.
\end{corollary}

The target grows: $2\theta'(\tau)=\log\frac\tau{2\pi}+O(\tau^{-2})$. At $\tau=3$
and $\Delta=\frac14$ the delay is $4.09\times10^{-8}$; by $\tau=10$ it is
$3\times10^{-27}$ and the target is $+0.464$.

\subsection{No correlator can wind}

The zero of Corollary~\ref{cor:protected} is an accident of half-integer
dimensions. The following is not.

Radial quantization of the ray about the pinned endpoint gives the
boundary-channel decomposition
$G(u)=\sum_n|\langle\Omega|O|n\rangle|^2e^{-\Delta_nu}=\int_{[0,\infty)}e^{-\Delta
u}\dd\mu(\Delta)$ with $\mu\geq0$, using only that the defect Hilbert space has a
positive inner product and that the dilatation generator is bounded below. For
$G_\Delta$ this is the primary with its descendants,
$\mu=\sum_k\frac{\Gamma(2\Delta+k)}{k!\Gamma(2\Delta)}\delta_{\Delta+k}$, with
manifestly positive coefficients. By Bernstein's theorem $G$ is then completely
monotone and $\widehat k_+$ is a Stieltjes transform.

\begin{proposition}[No winding]\label{prop:nowinding}
Let $G$ be the ray kernel of a two-point function in a unitary defect theory whose
dilatation generator is bounded below. Then $\re\widehat k_+>0$ on the whole real
axis --- so $\widehat k_+$ has no zeros there --- and for $\tau>0$,
$\arg\widehat k_+(\tau)\in(-\frac\pi2,0)$, whence
\[
 \int_0^T\mathcal T(\tau)\dd\tau<\frac\pi2\qquad\text{for every }T .
\]
If a real subtraction is needed for convergence the bound is $\pi$. Either way it
is a constant, while $\int_0^T2\theta'=2\theta(T)$ is $175.9$ at $T=100$ and
$4069.1$ at $T=10^3$.
\end{proposition}

\begin{proof}
$\widehat k_+(\tau)=\int(\Delta+i\tau)^{-1}\dd\mu$ has
$\re\widehat k_+=\int\Delta(\Delta^2+\tau^2)^{-1}\dd\mu>0$ and
$\im\widehat k_+=-\tau\int(\Delta^2+\tau^2)^{-1}\dd\mu<0$. Then
$\int_0^T\mathcal T=\arg\widehat k_+(0^+)-\arg\widehat k_+(T)$.
\end{proof}

Two readings. \textbf{Phase winding is zeros crossing the contour}, and a positive
spectral measure forbids them: the scattering dictionary of
Section~\ref{sec:allpass} lists the zeros of $\xi$ as the resonances of
$K_\omega$, and a correlator of a unitary theory has none to list. And
\textbf{the obstruction is that the generator is bounded below}, not that its
spectrum is discrete --- a continuum on $[0,\infty)$ gives the same Stieltjes
transform.

\begin{remark}[The three escapes, one of which is already a trap]\label{rem:escapes}
Given those hypotheses, winding requires exactly one of: a signed measure, that
is a non-unitary theory; a spectrum \emph{not} bounded below, so that the
imaginary axis lies outside the region of absolute convergence and the object on
it is a continuation past poles --- which is precisely
Remark~\ref{rem:trap}, seen from the candidate's side rather than the target's;
or a \emph{ratio}. The Weil family takes the third route, and
Proposition~\ref{prop:groupdelay} says so: the same real function
$b(\tau^2)+w_0$ is the archimedean multiplier of the form, where it is a positive
weight with no phase, and the phase derivative of the transfer, where it winds. A
correlator offers only the first reading.
\end{remark}

\subsection{What does wind, and why that is a problem}

The phase shift of the inverted harmonic oscillator on the half-line is
$\eta(\tau)=\arg\Gamma(\frac14+\frac{i\tau}2)$ \cite{BhaduriKhareLaw}, so
\begin{equation}\label{eq:freewinding}
 2\eta'(\tau)-2\theta'(\tau)=\log\pi\qquad\text{exactly, for every }\tau,
\end{equation}
verified to $10^{-14}$. That system is $\frac12(xp+px)$ --- the dilatation
generator --- on a ray, with continuous spectrum on all of $\R$, which is what
Proposition~\ref{prop:nowinding} requires. Two consequences, in opposite
directions. The delay clause is satisfied by a \emph{free} one-dimensional
Hamiltonian, so as a test of a theory it has almost no power. And what the free
generator misses is exactly the conductor: the constant $\log\pi$, which is the
$\pi^{-s/2}$ of the completion.

That the target is an inner function is the same statement:
$e^{2i\theta(\tau)}=\pi^{-i\tau}\Gamma(\frac14+\frac{i\tau}2)/
\Gamma(\frac14-\frac{i\tau}2)$, unimodular, with poles at $\tau=i(\frac12+2n)$ and
zeros at $\tau=-i(\frac12+2n)$. Proposition~\ref{prop:flux} already says
$K_\omega$ is inner; Proposition~\ref{prop:nowinding} says no correlator is.

\subsection{The expansion, and three graded invariants}

The discrimination is in the expansion. Let a candidate's reflection amplitude be
unimodular of $\Gamma$-ratio type,
\begin{equation}\label{eq:gammaratio}
 S(\tau)=e^{2ic_1\tau}\prod_j
 \left(\frac{\Gamma(a_j+i\beta_j\tau)}{\Gamma(a_j-i\beta_j\tau)}\right)^{n_j},
 \qquad a_j>0,\ \beta_j>0,\ n_j>0,
\end{equation}
finitely many factors. From $\re\psi(a+iy)=\log y+\frac{B_2(a)}{2y^2}+O(y^{-4})$
with $B_2(a)=a^2-a+\frac16$:

\begin{proposition}[The delay expansion]\label{prop:expansion}
\[
 \mathcal T(\tau)=\Lambda\log\tau+c_0
 +\frac1{\tau^2}\sum_j\frac{n_jB_2(a_j)}{\beta_j}+O(\tau^{-4}),
 \qquad \Lambda=2\sum_jn_j\beta_j,
\]
with $c_0=2c_1+2\sum_jn_j\beta_j\log\beta_j$. The target is
$\mathcal T_{\rm Weil}=\log\tau-\log2\pi-\frac1{24\tau^2}+O(\tau^{-4})$.
\end{proposition}

\begin{corollary}[$\Lambda$ is the degree]\label{cor:degree}
For the archimedean factor $\prod_{j=1}^d\Gamma_\R(s+\mu_j)$ of a degree-$d$
$L$-function at $s=\frac12+i\tau$ one has $a_j=\frac{1/2+\mu_j}2$,
$\beta_j=\frac12$, $n_j=1$, so $\Lambda=d$ exactly. Matching the target says:
degree one.
\end{corollary}

\begin{corollary}[The shift must sit near the critical line]\label{cor:signrule}
Matching the $\tau^{-2}$ coefficient requires
$\sum_jn_jB_2(a_j)/\beta_j=-\frac1{24}<0$, hence at least one $a_j$ with
$B_2(a_j)<0$, that is
\[
 \Big|a_j-\tfrac12\Big|<\tfrac1{2\sqrt3}=0.288675\ldots
\]
\end{corollary}

\begin{center}
\begin{tabular}{@{}lccc@{}}
\toprule
factor, at $s=\frac12+i\tau$ & $a$ & $B_2(a)$ & in the window\\
\midrule
$\Gamma_\R(s)$ --- $\zeta$, even $\chi$ & $\frac14$ & $-\frac1{48}$ & yes\\
$\Gamma_\R(s+1)$ --- odd $\chi$ & $\frac34$ & $-\frac1{48}$ & yes\\
$\Gamma_\C(s)$ --- a complex place & $\frac12$ & $-\frac1{12}$ & yes, deepest\\
$\Gamma(1+i\,\cdot)$ --- generic reflection amplitude & $1$ & $+\frac16$ & \textbf{no}\\
$\Gamma(\frac32+i\,\cdot)$ & $\frac32$ & $+\frac{11}{12}$ & \textbf{no}\\
\bottomrule
\end{tabular}
\end{center}

\begin{proposition}[Rigidity at a common shift]\label{prop:rigidity}
If all $a_j=a$, then $\Lambda=1$ together with
$\sum_jn_jB_2(a)/\beta_j=-\frac1{24}$ forces
\[
 \sum_jn_j\ \leq\ \big(48|B_2(a)|\big)^{-1/2},
\]
with equality if and only if all $\beta_j$ agree. The bound equals one exactly
when $a\in\{\frac14,\frac34\}$, so with integer multiplicities there is a single
factor, $n=1$, $\beta=\frac12$, $a\in\{\frac14,\frac34\}$: the archimedean factor
of a degree-one $L$-function, of either parity, and nothing else.
\end{proposition}

\begin{proof}
The two conditions read $\sum_jn_j\beta_j=\frac12$ and
$\sum_jn_j/\beta_j=(24|B_2(a)|)^{-1}$. Cauchy--Schwarz gives
\[
 \Big(\sum_jn_j\Big)^2\leq\Big(\sum_jn_j\beta_j\Big)\Big(\sum_jn_j/\beta_j\Big),
\]
with equality if and only if the $\beta_j$ agree; and $|B_2(a)|=\frac1{48}$ has
roots $a=\frac14$ and $a=\frac34$.
\end{proof}

The common-shift hypothesis is not removable: with $a_1=\frac12$, $a_2=1$,
$n_1=n_2=1$, $\beta_1=0.1576708$, $\beta_2=0.3423292$ both conditions hold, a
positive $B_2$ paying for a negative one. What is unconditional is
Corollary~\ref{cor:signrule}.

\subsection{Liouville, excluded}

The bulk Liouville reflection amplitude is unimodular and does wind:
$\frac{d}{dP}\arg S=4Q\log P+C_L+\frac{Q}{12P^2}+O(P^{-4})$, $Q=b+b^{-1}$.
The Liouville coordinate is the radial one, so the momenta are linearly related,
$\tau=2\kappa P$, and matching the leading coefficient \emph{forces} $\kappa=2Q$.
Nothing is then left but $b$; the cosmological constant $\mu$ enters the
prefactor linearly in $P$ and so moves only $c_0$. Substituting $P=\tau/4Q$,
\[
 \mathcal T_L(\tau)=\log\tau-\log4Q+\frac{C_L}{4Q}+\frac{Q^2}{3\tau^2}+O(\tau^{-4}) .
\]
The constant can be matched, by one $\mu$ for each $b$. The $\tau^{-2}$ term
cannot: $Q\geq2$ for every $b>0$, so $Q^2/3\geq\frac43$, against the target's
$-\frac1{24}$ --- wrong sign, and never within a factor of $32$ in magnitude.
This is Corollary~\ref{cor:signrule} with $a=1$. \textbf{Liouville is excluded at
every coupling}, and so is every amplitude built from $\Gamma(m+i\,\cdot)$ with
$m\geq1$.

\begin{remark}[Where the primes are not]\label{rem:notprimes}
None of this expansion sees the primes. They are in the bounded fluctuation
$S(\tau)=\pi^{-1}\arg\zeta(\frac12+i\tau)$, which is invisible to every
asymptotic order. So the expansion tests whether a candidate has the right
\emph{archimedean} data, to any depth wanted, and says nothing about whether it
has any arithmetic --- two questions that were previously entangled in one test.
This is the same separation Section~\ref{sec:sampling} found on the other side,
where the margin turned out to be set by the density and not by the primes.
\end{remark}

\begin{remark}[Status of this section]\label{rem:delaystatus}
Propositions~\ref{prop:defectdelay}, \ref{prop:nowinding}, \ref{prop:expansion}
and \ref{prop:rigidity} and their corollaries are written proofs, unconditional;
Proposition~\ref{prop:nowinding} quotes Bernstein's theorem. The Liouville
exclusion assumes the momentum dictionary is linear, which is stated where it is
used and which Section~\ref{sec:delay}'s other results do not need. Every number
is verified in Appendix~\ref{app:checks}. Section 11.1's description of
\cite{OWL} is a reading with its scope in the footnote; the defect spectrum and
normalizations are quoted, not re-derived.
\end{remark}
